% STATUS: written 2026-10-12
% STATUS: extended 2026-11-07 (Idrissi and Oukhtite 2019 read in full from the
% library PDF)
% Sources read in full: Herstein 1982 (pp. 356--357, the introduction and Theorem 1), Lanski 1977 (pp. 449--460), Lanski 1985 (p. 765, the abstract and introduction), Martindale and Miers 1986 (p. 14);
% Idrissi and Oukhtite 2019 (pp. 1950001-1 to 1950001-11).
\section{Rings with involution and *-derivations}\label{sec:prime-semiprime:involution}

Section~\ref{sec:prime-semiprime:ideals-lie-ideals} ended with the skew elements
of a ring with involution. An \emph{involution} is an anti-automorphism of period
at most $2$ \cite[p.~14]{MartindaleMiers1986}. The symmetric elements $s^*=s$ form
a Jordan subring, and the skew elements $k^*=-k$ form a Lie subring
\cite[p.~14]{MartindaleMiers1986}. We ask what a derivation does on these sets.

Herstein took the commuting condition of
Section~\ref{sec:prime-semiprime:herstein-school} to this setting
\cite[p.~356]{Herstein1982}. He asks that $d(s)d(t)=d(t)d(s)$ only for symmetric
$s$ and $t$ \cite[p.~356]{Herstein1982}.

The full result of the earlier section cannot hold here \cite[p.~356]{Herstein1982}.

\begin{example}[{\cite[p.~356]{Herstein1982}}]\label{ex:prime:herstein-symplectic}
Let $R$ be the ring of $2\times 2$ matrices over a field. For the symplectic
involution, all symmetric elements are central, so the values of any derivation
on symmetric elements commute. For the transpose, take
\[
d(x)=xe_{11}-e_{11}x .
\]
Then $d(x)d(y)=d(y)d(x)$ for symmetric $x$ and $y$.
\end{example}

So Herstein proves a weaker
conclusion.

\begin{theorem}[{\cite[Theorem~2, p.~366]{Herstein1982}}]\label{thm:prime:herstein-involution}
Let $R$ be a prime ring with involution, of characteristic not $2$, and let
$d\neq 0$ be a derivation of $R$ with $d(s)d(t)=d(t)d(s)$ for all symmetric $s$
and $t$. Then $R$ is commutative, or $R$ is an order in a simple algebra of
dimension four.
\end{theorem}

Lanski studied Lie ideals and derivations in rings with involution
\cite[p.~449]{Lanski1977}.
Let $L$ be a Lie ideal. Suppose that the commutator of an element with $L$ is
central. Then the element or $L$ is central, unless the ring is an order in a
simple algebra of dimension at most four \cite[p.~449]{Lanski1977}. This holds for
the symmetric part or the skew part of $L$ \cite[p.~449]{Lanski1977}.

Lanski applies this to derivations \cite[p.~449]{Lanski1977}. Suppose a
derivation sends the symmetric or skew part of $L$ into the center. Then the same
three cases occur \cite[p.~449]{Lanski1977}.

\begin{theorem}[{\cite[Theorem~9, p.~459]{Lanski1977}}]\label{thm:prime:lanski-involution}
Let $R$ be a prime ring of characteristic not $2$ with an involution, let $D$ be
a derivation of $R$, and let $L$ be a Lie ideal with $L^*=L$ such that $D(L)$ lies
in the skew elements or in the symmetric elements. Then $L\subseteq Z$, or
$D(R)=0$, or $R$ satisfies the standard identity of degree four.
\end{theorem}

Section~\ref{sec:prime-semiprime:kharchenko} has a version with involution. Every
generalized identity with derivations and the involution follows from the
generalized identities with involution that $R$ satisfies
\cite[p.~765]{Lanski1985}. This gives a general setting for the result of
Herstein \cite[p.~765]{Lanski1985}.

In characteristic zero, no outer derivation of $R$ acts algebraically on the
symmetric elements \cite[p.~765]{Lanski1985}. Suppose the images of the
symmetric elements under the derivations satisfy a polynomial relation. Then
$R$ satisfies a generalized polynomial identity \cite[p.~765]{Lanski1985}.

An involution is of the \emph{second kind} when its symmetric elements meet the
center in a nonzero set \cite[p.~1950001-1]{Idrissi2019}. Idrissi and Oukhtite
take the commuting condition of
Section~\ref{sec:prime-semiprime:herstein-school} and put the commutator of an
element with its image in place of the element itself
\cite[p.~1950001-4]{Idrissi2019}.

\begin{theorem}[{\cite[Theorem~1, p.~1950001-4]{Idrissi2019}}]\label{thm:prime:idrissi-commutator}
Let $(R,*)$ be a $2$-torsion free prime ring with an involution of the second kind,
and let $F$ be a nonzero generalized derivation of $R$ associated with a derivation
$d$. Then $R$ is commutative if and only if $F([x,x^{*}])\in Z(R)$ for all $x\in R$.
\end{theorem}

They argue the same way for the Jordan product, and again the ring has to commute
\cite[p.~1950001-6]{Idrissi2019}.

\begin{theorem}[{\cite[Theorem~2, p.~1950001-6]{Idrissi2019}}]\label{thm:prime:idrissi-jordan}
Let $(R,*)$ be a $2$-torsion free prime ring with an involution of the second kind.
If $R$ admits a nonzero generalized derivation $F$ associated with a derivation $d$,
then $R$ is commutative if and only if $F(x\circ x^{*})\in Z(R)$ for all $x\in R$.
\end{theorem}

A corollary replaces the condition $F([x,x^{*}])\in Z(R)$ by
$F([x,y])\pm[x,y]\in Z(R)$ for all $x$ and $y$, and adds that the three statements
are equivalent \cite[Corollary~1, p.~1950001-5]{Idrissi2019}. The authors also note
that an earlier result for a square closed Lie ideal fails here. When the ideal is
the whole ring, that result would force the ring to be commutative
\cite[p.~1950001-5]{Idrissi2019}.

Section~\ref{sec:prime-semiprime:jordan-ideals} turns from Lie ideals to Jordan
ideals.

Nonlinearity does not weaken the conclusion. On a prime ring with involution
whose skew elements are $K$, a map $f$ with $[f(x),f(y)]=[x,y]$ for all
$x,y\in K$ is additive and preserves the commutator
\cite[Abstract]{NonlinearStrongCommuting2012}.

For two generalized derivations on a $*$-ideal of a noncommutative prime $*$-ring, the $*$-identities they satisfy force the same conclusion Herstein's theorem draws~\cite{WlHersteinsTheorem2023}.

For two generalized derivations on a $*$-ideal of a noncommutative prime $*$-ring, the $*$-identities they satisfy force the conclusion Herstein's theorem draws~\cite{WlHersteinSTheoremForPairOfGener2023}.

Swain's determination of the Lie derivations of the skew elements of a prime ring with involution, which is the Lie form of the skew-element question the chapter treats~\cite{WlLieDerivationsOfTheSkewElement1996}.
